% Part: computability
% Chapter: computability-theory
% Section: computable-sets

\documentclass[../../../include/open-logic-section]{subfiles}

\begin{document}

\olfileid{cmp}{thy}{cps}
\olsection{المجموعات القابلة للحساب}

وننقل الآن مفهوم قابلية الحساب من الدوال إلى
المجموعات، على الوجه الآتي:

\begin{defn}
  إذا أخذنا مجموعة $S$ من الأعداد الطبيعية، قلنا إن $S$ \emph{قابلة للحساب}
  متى كانت دالتها المميزة~$\Char{S}$ قابلة للحساب، وبالعكس. فشرط
  قابلية $S$ للحساب هو قابلية الدالة
\[
\Char{S}(x) =
\begin{cases}
1 & \text{إذا كان $x \in S$} \\
0 & \text{في غير ذلك}
\end{cases}
\]
للحساب. وعلى هذا تُعد العلاقة
$R(x_0,\dots,x_{k-1})$ قابلة للحساب إذا وفقط إذا كانت دالتها المميزة
كذلك.

وتسمى المجموعات والعلائق القابلة للحساب أيضًا \emph{قابلة للقرار}.
\end{defn}

\begin{explain}
فقد تنوع الآن متعلق قابلية الحساب: تارة الدوال الجزئية، وتارة الدوال،
وتارة المجموعات. ولا بد من التمييز بينها.{} \iftag{TMs}{فآلة تورنغ الحاسبة
لدالة جزئية تعطي قيمة الدالة عند كل دخل تُعرّف فيه؛ وأما
الحاسبة لمجموعة فتعطي $1$ أو~$0$ بحسب حكمها
بانتماء الدخل إلى المجموعة أو عدم انتمائه.}{}
\end{explain}

\end{document}
